\documentclass[12pt]{article}
\usepackage[margin=1in, top=0.75in, bottom=0.75in]{geometry}
% worksheet-preamble.tex — the shipped preamble every generated document
% \input's. SOURCE OF TRUTH: this file. references/latex-templates.md documents
% these macros but no longer serves as the text agents transcribe — retyping a
% ~130-line preamble per document, three documents per run, is exactly the
% transcription-drift error class the rest of the pipeline exists to catch.
%
% Usage (compile.sh stages this file beside your .tex automatically):
%   \documentclass[12pt]{article}
%   \usepackage[margin=1in, top=0.75in, bottom=0.75in]{geometry}
%   % worksheet-preamble.tex — the shipped preamble every generated document
% \input's. SOURCE OF TRUTH: this file. references/latex-templates.md documents
% these macros but no longer serves as the text agents transcribe — retyping a
% ~130-line preamble per document, three documents per run, is exactly the
% transcription-drift error class the rest of the pipeline exists to catch.
%
% Usage (compile.sh stages this file beside your .tex automatically):
%   \documentclass[12pt]{article}
%   \usepackage[margin=1in, top=0.75in, bottom=0.75in]{geometry}
%   % worksheet-preamble.tex — the shipped preamble every generated document
% \input's. SOURCE OF TRUTH: this file. references/latex-templates.md documents
% these macros but no longer serves as the text agents transcribe — retyping a
% ~130-line preamble per document, three documents per run, is exactly the
% transcription-drift error class the rest of the pipeline exists to catch.
%
% Usage (compile.sh stages this file beside your .tex automatically):
%   \documentclass[12pt]{article}
%   \usepackage[margin=1in, top=0.75in, bottom=0.75in]{geometry}
%   \input{worksheet-preamble}
%   \input{figure-macros}          % when using the shipped figure macros
%
% geometry stays in the document, NOT here: worksheets use margin=1in
% (top/bottom 0.75in) while study guides use margin=0.85in (top/bottom 0.7in),
% and geometry options cannot be re-issued once the package is loaded.
%
% ASCII only in this file: pdflatex (the fallback engine) cannot typeset
% literal Unicode symbols, so documents must compile identically under both
% tectonic and pdflatex.

% ── Required packages (superset of the worksheet and study-guide shells; a
% worksheet loading mdframed it does not use costs nothing, one shared file
% that cannot drift is worth it) ─────────────────────────────────────────────
\usepackage{amsmath, amssymb}
\usepackage{tikz, pgfplots}
\usepackage{enumitem, fancyhdr, multicol, array, booktabs}
\usepackage{mdframed, xcolor}
\pgfplotsset{compat=1.18}
\usetikzlibrary{calc, angles, quotes}   % needed by the geometric figure macros

% ── Problem numbering ───────────────────────────────────────────────────────
\newcounter{prob}
\newcommand{\problem}[1]{\stepcounter{prob}\vspace{0.4cm}\noindent\textbf{\large\theprob.}\quad #1\vspace{0.3cm}}

% ── Shrink-to-fit title: \LARGE when it fits on one line, otherwise scale
% down to the text width — never enlarges, never wraps mid-phrase. (\resizebox
% comes from graphicx, which tikz already loads.) ────────────────────────────
\newlength{\fittedw}
\newcommand{\fittedtitle}[1]{%
  \settowidth{\fittedw}{\LARGE\bfseries #1}%
  \ifdim\fittedw>\linewidth\resizebox{\linewidth}{!}{\bfseries #1}%
  \else{\LARGE\bfseries #1}\fi}

% ── Running headers. Keep the left title SHORT (<= ~28 chars, e.g. "Triangle
% Trig Practice", not the full course name) — it shares the header line with
% the Name/Date blanks and a long title overlaps them. ──────────────────────
% \wsheader{Short Title} — student worksheet
\newcommand{\wsheader}[1]{%
  \pagestyle{fancy}\fancyhf{}%
  \fancyhead[L]{\textbf{\small #1}}%
  \fancyhead[R]{\small Name: \underline{\hspace{4.5cm}}~Date: \underline{\hspace{1.8cm}}}%
  \fancyfoot[C]{\thepage}%
  \renewcommand{\headrulewidth}{0.4pt}}

% \akheader{Short Topic} — answer key
\newcommand{\akheader}[1]{%
  \pagestyle{fancy}\fancyhf{}%
  \fancyhead[L]{\textbf{\small #1 --- Answer Key}}%
  \fancyhead[R]{\textit{\small For instructor/parent use}}%
  \fancyfoot[C]{\thepage}%
  \renewcommand{\headrulewidth}{0.4pt}}

% \ssheader{Topic} — skills summary / study guide
\newcommand{\ssheader}[1]{%
  \pagestyle{fancy}\fancyhf{}%
  \fancyhead[L]{\textbf{Skills Summary: #1}}%
  \fancyhead[R]{\small\textit{Study Guide \& Reference}}%
  \fancyfoot[C]{\thepage}%
  \renewcommand{\headrulewidth}{0.4pt}}

% ── Title blocks ────────────────────────────────────────────────────────────
% \wstitleblock{Title}{Course}{Date}
\newcommand{\wstitleblock}[3]{%
  \begin{center}
    \fittedtitle{#1}\\[0.3cm]
    {\large #2 \quad $\bullet$ \quad #3}
  \end{center}
  \noindent\rule{\linewidth}{0.4pt}\vspace{0.2cm}}

% \aktitleblock{Topic}{Course}{Date} — "Answer Key" is a subtitle BENEATH the
% topic, never appended to it: a long topic would wrap mid-phrase.
\newcommand{\aktitleblock}[3]{%
  \begin{center}
    \fittedtitle{#1}\\[0.15cm]
    {\Large\bfseries Answer Key}\\[0.25cm]
    {\large #2 \quad $\bullet$ \quad #3}
  \end{center}}

% \sstitleblock{Topic}
\newcommand{\sstitleblock}[1]{%
  \begin{center}
    {\LARGE\textbf{Skills Summary}}\\[0.2cm]
    {\large\textbf{#1}}\\[0.1cm]
    {\small\textit{Use this reference while completing your worksheet or when studying.}}
  \end{center}
  \vspace{0.1cm}
  \noindent\rule{\linewidth}{1.5pt}
  \vspace{0.3cm}}

% ── Study-guide color palette and boxes ─────────────────────────────────────
\definecolor{skillblue}{RGB}{30,100,180}
\definecolor{skillbluebg}{RGB}{235,244,255}
\definecolor{warnorange}{RGB}{200,90,0}
\definecolor{warnbg}{RGB}{255,243,230}
\definecolor{exgreen}{RGB}{20,120,60}
\definecolor{exgreenbg}{RGB}{230,248,238}

% Formula/rule box
\newmdenv[
  backgroundcolor=skillbluebg,
  linecolor=skillblue, linewidth=1.5pt,
  innertopmargin=6pt, innerbottommargin=6pt,
  innerleftmargin=10pt, innerrightmargin=10pt,
  skipabove=6pt, skipbelow=4pt
]{formulabox}

% Mini example box
\newmdenv[
  backgroundcolor=exgreenbg,
  linecolor=exgreen, linewidth=1pt,
  innertopmargin=5pt, innerbottommargin=5pt,
  innerleftmargin=10pt, innerrightmargin=10pt,
  skipabove=4pt, skipbelow=4pt
]{examplebox}

% Watch-out box
\newmdenv[
  backgroundcolor=warnbg,
  linecolor=warnorange, linewidth=1pt,
  innertopmargin=5pt, innerbottommargin=5pt,
  innerleftmargin=10pt, innerrightmargin=10pt,
  skipabove=4pt, skipbelow=4pt
]{watchoutbox}

% Skill section heading
\newcommand{\skillheading}[1]{%
  \vspace{0.4cm}
  {\large\textbf{\textcolor{skillblue}{#1}}}
  \vspace{0.1cm}
  \hrule height 1pt
  \vspace{0.2cm}
}

%   % figure-macros.tex — shipped geometric figure macros. SOURCE OF TRUTH: this
% file; references/latex-templates.md documents usage. Requires the tikz
% libraries loaded by worksheet-preamble.tex (calc, angles, quotes), so
% \input worksheet-preamble first.
%
% CONTRACT with the layout/prose checkers (do not break it):
%   - Macro names contain rt/tri/fig — that is how tests/check_layout.py and
%     tests/check_prose_consistency.py recognize a macro-style figure.
%   - Mandatory arguments are FLAT braces only (no nested {} inside an arg);
%     the checkers' arg scanners match [^{}]* per group.
%   - Value-bearing args carry exactly the printed label text; scale/styling
%     goes in the single optional [..] argument, which the checkers ignore
%     (so scale=1.2 is never mistaken for a figure value).
%
% Figure conventions (from SKILL.md / latex-templates.md): every number
% printed in a figure must come from the problem's verify JSON; vertices are
% A/B/C with sides a/b/c OPPOSITE them — the same convention verify.py's
% triangle type uses. Right triangles put the right angle at C, so c (= AB)
% is the hypotenuse — the SAME convention scripts/render_figures.py constructs
% (its RIGHT_ANGLE_VERTEX constant): \refrt sits beside renderer-built
% \probfig figures on the page, so the two sources must never disagree.
% tests/test_figure_convention.py enforces this by parsing the
% "% right angle mark" lines below — that trailing comment is load-bearing;
% keep it on every right-angle-mark \draw.
%
% ASCII only: pdflatex (the fallback engine) cannot typeset literal Unicode.

% ── \rtfig[scale]{bottom-leg label}{right-leg label}{hypotenuse label} ──────
% Right triangle, right angle at C (hypotenuse c = AB — the renderer's
% convention). A=(0,0), C=(4,0), B=(4,3): the 3-4-5 shape keeps every label
% clear of the drawn lines at default scale. Bottom leg = b (AC), right leg
% = a (CB), hypotenuse = c (AB) — name sides accordingly in the labels.
% Example: \rtfig{$8$}{$6$}{$x$}   or   \rtfig[0.9]{$b = 8$}{$a = 6$}{$c$}
\newcommand{\rtfig}[4][1.2]{%
  \begin{center}
  \begin{tikzpicture}[scale=#1]
    \coordinate (A) at (0,0);
    \coordinate (B) at (4,3);
    \coordinate (C) at (4,0);
    \draw[thick] (A) -- (C) -- (B) -- cycle;
    \draw (C) ++(-.25,0) -- ++(0,.25) -- ++(.25,0);  % right angle mark
    \node[below left] at (A) {$A$};
    \node[above right] at (B) {$B$};
    \node[below right] at (C) {$C$};
    \node[below] at (2,0) {#2};
    \node[right] at (4,1.5) {#3};
    \node[above left] at (2,1.5) {#4};
  \end{tikzpicture}
  \end{center}}

% ── \trifig[scale]{c}{b}{A-deg}{c-label}{b-label}{a-label}{angle-label} ─────
% General triangle, TO SCALE by construction (SAS): A at the origin, B on the
% x-axis at distance c, C computed from side b and angle A (TikZ trig takes
% degrees). Numeric args #2-#4 are the ACTUAL given values, so the drawing
% cannot disagree with the math; label args #5-#8 carry what is printed.
% Example: \trifig{7}{5}{34}{$c = 7$}{$b = 5$}{$a = ?$}{$34^\circ$}
\newcommand{\trifig}[8][0.6]{%
  \begin{center}
  \begin{tikzpicture}[scale=#1]
    \coordinate (A) at (0,0);
    \coordinate (B) at (#2,0);
    \coordinate (C) at ({#3*cos(#4)},{#3*sin(#4)});
    \draw[thick] (A) -- (B) -- (C) -- cycle;
    \node[below left]  at (A) {$A$};
    \node[below right] at (B) {$B$};
    \node[above]       at (C) {$C$};
    \node[below]       at ($(A)!0.5!(B)$) {#5};
    \node[above left]  at ($(A)!0.5!(C)$) {#6};
    \node[above right] at ($(B)!0.5!(C)$) {#7};
    \pic [draw, angle radius=7mm, angle eccentricity=1.6, "#8"] {angle = B--A--C};
  \end{tikzpicture}
  \end{center}}

% ── \refrt — the value-free reference figure ────────────────────────────────
% SKILL.md's figure-scope rule requires shared labelling conventions to live
% in ONE value-free reference figure placed with the directions; until now no
% template for it existed, so every generation invented one — and an invented
% figure that accidentally carried a number would itself trip the
% all-or-nothing figure rule it exists to satisfy. This macro has zero
% numerals in its output by construction (it takes no arguments and prints
% only letters), and the caption is baked in verbatim so the figure cannot be
% mistaken for a problem's givens. Right angle at C, so sides read: a = BC
% (opposite A), b = AC (opposite B), c = AB (opposite C, the hypotenuse) —
% matching render_figures.py's RIGHT_ANGLE_VERTEX, because this reference
% figure defines the convention students apply to every renderer-built
% figure on the sheet.
\newcommand{\refrt}{%
  \begin{center}
  \begin{tikzpicture}[scale=1.0]
    \coordinate (A) at (0,0);
    \coordinate (B) at (3.6,2.7);
    \coordinate (C) at (3.6,0);
    \draw[thick] (A) -- (C) -- (B) -- cycle;
    \draw (C) ++(-.25,0) -- ++(0,.25) -- ++(.25,0);  % right angle mark
    \node[below left] at (A) {$A$};
    \node[above right] at (B) {$B$};
    \node[below right] at (C) {$C$};
    \node[below] at (1.8,0) {$b$};
    \node[right] at (3.6,1.35) {$a$};
    \node[above left] at (1.8,1.35) {$c$};
    \pic [draw, angle radius=6mm, angle eccentricity=1.7, "$A$"] {angle = C--A--B};
  \end{tikzpicture}\\
  {\small\textit{How every triangle here is labelled. No values shown: use the
  numbers given in each problem.}}
  \end{center}}
          % when using the shipped figure macros
%
% geometry stays in the document, NOT here: worksheets use margin=1in
% (top/bottom 0.75in) while study guides use margin=0.85in (top/bottom 0.7in),
% and geometry options cannot be re-issued once the package is loaded.
%
% ASCII only in this file: pdflatex (the fallback engine) cannot typeset
% literal Unicode symbols, so documents must compile identically under both
% tectonic and pdflatex.

% ── Required packages (superset of the worksheet and study-guide shells; a
% worksheet loading mdframed it does not use costs nothing, one shared file
% that cannot drift is worth it) ─────────────────────────────────────────────
\usepackage{amsmath, amssymb}
\usepackage{tikz, pgfplots}
\usepackage{enumitem, fancyhdr, multicol, array, booktabs}
\usepackage{mdframed, xcolor}
\pgfplotsset{compat=1.18}
\usetikzlibrary{calc, angles, quotes}   % needed by the geometric figure macros

% ── Problem numbering ───────────────────────────────────────────────────────
\newcounter{prob}
\newcommand{\problem}[1]{\stepcounter{prob}\vspace{0.4cm}\noindent\textbf{\large\theprob.}\quad #1\vspace{0.3cm}}

% ── Shrink-to-fit title: \LARGE when it fits on one line, otherwise scale
% down to the text width — never enlarges, never wraps mid-phrase. (\resizebox
% comes from graphicx, which tikz already loads.) ────────────────────────────
\newlength{\fittedw}
\newcommand{\fittedtitle}[1]{%
  \settowidth{\fittedw}{\LARGE\bfseries #1}%
  \ifdim\fittedw>\linewidth\resizebox{\linewidth}{!}{\bfseries #1}%
  \else{\LARGE\bfseries #1}\fi}

% ── Running headers. Keep the left title SHORT (<= ~28 chars, e.g. "Triangle
% Trig Practice", not the full course name) — it shares the header line with
% the Name/Date blanks and a long title overlaps them. ──────────────────────
% \wsheader{Short Title} — student worksheet
\newcommand{\wsheader}[1]{%
  \pagestyle{fancy}\fancyhf{}%
  \fancyhead[L]{\textbf{\small #1}}%
  \fancyhead[R]{\small Name: \underline{\hspace{4.5cm}}~Date: \underline{\hspace{1.8cm}}}%
  \fancyfoot[C]{\thepage}%
  \renewcommand{\headrulewidth}{0.4pt}}

% \akheader{Short Topic} — answer key
\newcommand{\akheader}[1]{%
  \pagestyle{fancy}\fancyhf{}%
  \fancyhead[L]{\textbf{\small #1 --- Answer Key}}%
  \fancyhead[R]{\textit{\small For instructor/parent use}}%
  \fancyfoot[C]{\thepage}%
  \renewcommand{\headrulewidth}{0.4pt}}

% \ssheader{Topic} — skills summary / study guide
\newcommand{\ssheader}[1]{%
  \pagestyle{fancy}\fancyhf{}%
  \fancyhead[L]{\textbf{Skills Summary: #1}}%
  \fancyhead[R]{\small\textit{Study Guide \& Reference}}%
  \fancyfoot[C]{\thepage}%
  \renewcommand{\headrulewidth}{0.4pt}}

% ── Title blocks ────────────────────────────────────────────────────────────
% \wstitleblock{Title}{Course}{Date}
\newcommand{\wstitleblock}[3]{%
  \begin{center}
    \fittedtitle{#1}\\[0.3cm]
    {\large #2 \quad $\bullet$ \quad #3}
  \end{center}
  \noindent\rule{\linewidth}{0.4pt}\vspace{0.2cm}}

% \aktitleblock{Topic}{Course}{Date} — "Answer Key" is a subtitle BENEATH the
% topic, never appended to it: a long topic would wrap mid-phrase.
\newcommand{\aktitleblock}[3]{%
  \begin{center}
    \fittedtitle{#1}\\[0.15cm]
    {\Large\bfseries Answer Key}\\[0.25cm]
    {\large #2 \quad $\bullet$ \quad #3}
  \end{center}}

% \sstitleblock{Topic}
\newcommand{\sstitleblock}[1]{%
  \begin{center}
    {\LARGE\textbf{Skills Summary}}\\[0.2cm]
    {\large\textbf{#1}}\\[0.1cm]
    {\small\textit{Use this reference while completing your worksheet or when studying.}}
  \end{center}
  \vspace{0.1cm}
  \noindent\rule{\linewidth}{1.5pt}
  \vspace{0.3cm}}

% ── Study-guide color palette and boxes ─────────────────────────────────────
\definecolor{skillblue}{RGB}{30,100,180}
\definecolor{skillbluebg}{RGB}{235,244,255}
\definecolor{warnorange}{RGB}{200,90,0}
\definecolor{warnbg}{RGB}{255,243,230}
\definecolor{exgreen}{RGB}{20,120,60}
\definecolor{exgreenbg}{RGB}{230,248,238}

% Formula/rule box
\newmdenv[
  backgroundcolor=skillbluebg,
  linecolor=skillblue, linewidth=1.5pt,
  innertopmargin=6pt, innerbottommargin=6pt,
  innerleftmargin=10pt, innerrightmargin=10pt,
  skipabove=6pt, skipbelow=4pt
]{formulabox}

% Mini example box
\newmdenv[
  backgroundcolor=exgreenbg,
  linecolor=exgreen, linewidth=1pt,
  innertopmargin=5pt, innerbottommargin=5pt,
  innerleftmargin=10pt, innerrightmargin=10pt,
  skipabove=4pt, skipbelow=4pt
]{examplebox}

% Watch-out box
\newmdenv[
  backgroundcolor=warnbg,
  linecolor=warnorange, linewidth=1pt,
  innertopmargin=5pt, innerbottommargin=5pt,
  innerleftmargin=10pt, innerrightmargin=10pt,
  skipabove=4pt, skipbelow=4pt
]{watchoutbox}

% Skill section heading
\newcommand{\skillheading}[1]{%
  \vspace{0.4cm}
  {\large\textbf{\textcolor{skillblue}{#1}}}
  \vspace{0.1cm}
  \hrule height 1pt
  \vspace{0.2cm}
}

%   % figure-macros.tex — shipped geometric figure macros. SOURCE OF TRUTH: this
% file; references/latex-templates.md documents usage. Requires the tikz
% libraries loaded by worksheet-preamble.tex (calc, angles, quotes), so
% \input worksheet-preamble first.
%
% CONTRACT with the layout/prose checkers (do not break it):
%   - Macro names contain rt/tri/fig — that is how tests/check_layout.py and
%     tests/check_prose_consistency.py recognize a macro-style figure.
%   - Mandatory arguments are FLAT braces only (no nested {} inside an arg);
%     the checkers' arg scanners match [^{}]* per group.
%   - Value-bearing args carry exactly the printed label text; scale/styling
%     goes in the single optional [..] argument, which the checkers ignore
%     (so scale=1.2 is never mistaken for a figure value).
%
% Figure conventions (from SKILL.md / latex-templates.md): every number
% printed in a figure must come from the problem's verify JSON; vertices are
% A/B/C with sides a/b/c OPPOSITE them — the same convention verify.py's
% triangle type uses. Right triangles put the right angle at C, so c (= AB)
% is the hypotenuse — the SAME convention scripts/render_figures.py constructs
% (its RIGHT_ANGLE_VERTEX constant): \refrt sits beside renderer-built
% \probfig figures on the page, so the two sources must never disagree.
% tests/test_figure_convention.py enforces this by parsing the
% "% right angle mark" lines below — that trailing comment is load-bearing;
% keep it on every right-angle-mark \draw.
%
% ASCII only: pdflatex (the fallback engine) cannot typeset literal Unicode.

% ── \rtfig[scale]{bottom-leg label}{right-leg label}{hypotenuse label} ──────
% Right triangle, right angle at C (hypotenuse c = AB — the renderer's
% convention). A=(0,0), C=(4,0), B=(4,3): the 3-4-5 shape keeps every label
% clear of the drawn lines at default scale. Bottom leg = b (AC), right leg
% = a (CB), hypotenuse = c (AB) — name sides accordingly in the labels.
% Example: \rtfig{$8$}{$6$}{$x$}   or   \rtfig[0.9]{$b = 8$}{$a = 6$}{$c$}
\newcommand{\rtfig}[4][1.2]{%
  \begin{center}
  \begin{tikzpicture}[scale=#1]
    \coordinate (A) at (0,0);
    \coordinate (B) at (4,3);
    \coordinate (C) at (4,0);
    \draw[thick] (A) -- (C) -- (B) -- cycle;
    \draw (C) ++(-.25,0) -- ++(0,.25) -- ++(.25,0);  % right angle mark
    \node[below left] at (A) {$A$};
    \node[above right] at (B) {$B$};
    \node[below right] at (C) {$C$};
    \node[below] at (2,0) {#2};
    \node[right] at (4,1.5) {#3};
    \node[above left] at (2,1.5) {#4};
  \end{tikzpicture}
  \end{center}}

% ── \trifig[scale]{c}{b}{A-deg}{c-label}{b-label}{a-label}{angle-label} ─────
% General triangle, TO SCALE by construction (SAS): A at the origin, B on the
% x-axis at distance c, C computed from side b and angle A (TikZ trig takes
% degrees). Numeric args #2-#4 are the ACTUAL given values, so the drawing
% cannot disagree with the math; label args #5-#8 carry what is printed.
% Example: \trifig{7}{5}{34}{$c = 7$}{$b = 5$}{$a = ?$}{$34^\circ$}
\newcommand{\trifig}[8][0.6]{%
  \begin{center}
  \begin{tikzpicture}[scale=#1]
    \coordinate (A) at (0,0);
    \coordinate (B) at (#2,0);
    \coordinate (C) at ({#3*cos(#4)},{#3*sin(#4)});
    \draw[thick] (A) -- (B) -- (C) -- cycle;
    \node[below left]  at (A) {$A$};
    \node[below right] at (B) {$B$};
    \node[above]       at (C) {$C$};
    \node[below]       at ($(A)!0.5!(B)$) {#5};
    \node[above left]  at ($(A)!0.5!(C)$) {#6};
    \node[above right] at ($(B)!0.5!(C)$) {#7};
    \pic [draw, angle radius=7mm, angle eccentricity=1.6, "#8"] {angle = B--A--C};
  \end{tikzpicture}
  \end{center}}

% ── \refrt — the value-free reference figure ────────────────────────────────
% SKILL.md's figure-scope rule requires shared labelling conventions to live
% in ONE value-free reference figure placed with the directions; until now no
% template for it existed, so every generation invented one — and an invented
% figure that accidentally carried a number would itself trip the
% all-or-nothing figure rule it exists to satisfy. This macro has zero
% numerals in its output by construction (it takes no arguments and prints
% only letters), and the caption is baked in verbatim so the figure cannot be
% mistaken for a problem's givens. Right angle at C, so sides read: a = BC
% (opposite A), b = AC (opposite B), c = AB (opposite C, the hypotenuse) —
% matching render_figures.py's RIGHT_ANGLE_VERTEX, because this reference
% figure defines the convention students apply to every renderer-built
% figure on the sheet.
\newcommand{\refrt}{%
  \begin{center}
  \begin{tikzpicture}[scale=1.0]
    \coordinate (A) at (0,0);
    \coordinate (B) at (3.6,2.7);
    \coordinate (C) at (3.6,0);
    \draw[thick] (A) -- (C) -- (B) -- cycle;
    \draw (C) ++(-.25,0) -- ++(0,.25) -- ++(.25,0);  % right angle mark
    \node[below left] at (A) {$A$};
    \node[above right] at (B) {$B$};
    \node[below right] at (C) {$C$};
    \node[below] at (1.8,0) {$b$};
    \node[right] at (3.6,1.35) {$a$};
    \node[above left] at (1.8,1.35) {$c$};
    \pic [draw, angle radius=6mm, angle eccentricity=1.7, "$A$"] {angle = C--A--B};
  \end{tikzpicture}\\
  {\small\textit{How every triangle here is labelled. No values shown: use the
  numbers given in each problem.}}
  \end{center}}
          % when using the shipped figure macros
%
% geometry stays in the document, NOT here: worksheets use margin=1in
% (top/bottom 0.75in) while study guides use margin=0.85in (top/bottom 0.7in),
% and geometry options cannot be re-issued once the package is loaded.
%
% ASCII only in this file: pdflatex (the fallback engine) cannot typeset
% literal Unicode symbols, so documents must compile identically under both
% tectonic and pdflatex.

% ── Required packages (superset of the worksheet and study-guide shells; a
% worksheet loading mdframed it does not use costs nothing, one shared file
% that cannot drift is worth it) ─────────────────────────────────────────────
\usepackage{amsmath, amssymb}
\usepackage{tikz, pgfplots}
\usepackage{enumitem, fancyhdr, multicol, array, booktabs}
\usepackage{mdframed, xcolor}
\pgfplotsset{compat=1.18}
\usetikzlibrary{calc, angles, quotes}   % needed by the geometric figure macros

% ── Problem numbering ───────────────────────────────────────────────────────
\newcounter{prob}
\newcommand{\problem}[1]{\stepcounter{prob}\vspace{0.4cm}\noindent\textbf{\large\theprob.}\quad #1\vspace{0.3cm}}

% ── Shrink-to-fit title: \LARGE when it fits on one line, otherwise scale
% down to the text width — never enlarges, never wraps mid-phrase. (\resizebox
% comes from graphicx, which tikz already loads.) ────────────────────────────
\newlength{\fittedw}
\newcommand{\fittedtitle}[1]{%
  \settowidth{\fittedw}{\LARGE\bfseries #1}%
  \ifdim\fittedw>\linewidth\resizebox{\linewidth}{!}{\bfseries #1}%
  \else{\LARGE\bfseries #1}\fi}

% ── Running headers. Keep the left title SHORT (<= ~28 chars, e.g. "Triangle
% Trig Practice", not the full course name) — it shares the header line with
% the Name/Date blanks and a long title overlaps them. ──────────────────────
% \wsheader{Short Title} — student worksheet
\newcommand{\wsheader}[1]{%
  \pagestyle{fancy}\fancyhf{}%
  \fancyhead[L]{\textbf{\small #1}}%
  \fancyhead[R]{\small Name: \underline{\hspace{4.5cm}}~Date: \underline{\hspace{1.8cm}}}%
  \fancyfoot[C]{\thepage}%
  \renewcommand{\headrulewidth}{0.4pt}}

% \akheader{Short Topic} — answer key
\newcommand{\akheader}[1]{%
  \pagestyle{fancy}\fancyhf{}%
  \fancyhead[L]{\textbf{\small #1 --- Answer Key}}%
  \fancyhead[R]{\textit{\small For instructor/parent use}}%
  \fancyfoot[C]{\thepage}%
  \renewcommand{\headrulewidth}{0.4pt}}

% \ssheader{Topic} — skills summary / study guide
\newcommand{\ssheader}[1]{%
  \pagestyle{fancy}\fancyhf{}%
  \fancyhead[L]{\textbf{Skills Summary: #1}}%
  \fancyhead[R]{\small\textit{Study Guide \& Reference}}%
  \fancyfoot[C]{\thepage}%
  \renewcommand{\headrulewidth}{0.4pt}}

% ── Title blocks ────────────────────────────────────────────────────────────
% \wstitleblock{Title}{Course}{Date}
\newcommand{\wstitleblock}[3]{%
  \begin{center}
    \fittedtitle{#1}\\[0.3cm]
    {\large #2 \quad $\bullet$ \quad #3}
  \end{center}
  \noindent\rule{\linewidth}{0.4pt}\vspace{0.2cm}}

% \aktitleblock{Topic}{Course}{Date} — "Answer Key" is a subtitle BENEATH the
% topic, never appended to it: a long topic would wrap mid-phrase.
\newcommand{\aktitleblock}[3]{%
  \begin{center}
    \fittedtitle{#1}\\[0.15cm]
    {\Large\bfseries Answer Key}\\[0.25cm]
    {\large #2 \quad $\bullet$ \quad #3}
  \end{center}}

% \sstitleblock{Topic}
\newcommand{\sstitleblock}[1]{%
  \begin{center}
    {\LARGE\textbf{Skills Summary}}\\[0.2cm]
    {\large\textbf{#1}}\\[0.1cm]
    {\small\textit{Use this reference while completing your worksheet or when studying.}}
  \end{center}
  \vspace{0.1cm}
  \noindent\rule{\linewidth}{1.5pt}
  \vspace{0.3cm}}

% ── Study-guide color palette and boxes ─────────────────────────────────────
\definecolor{skillblue}{RGB}{30,100,180}
\definecolor{skillbluebg}{RGB}{235,244,255}
\definecolor{warnorange}{RGB}{200,90,0}
\definecolor{warnbg}{RGB}{255,243,230}
\definecolor{exgreen}{RGB}{20,120,60}
\definecolor{exgreenbg}{RGB}{230,248,238}

% Formula/rule box
\newmdenv[
  backgroundcolor=skillbluebg,
  linecolor=skillblue, linewidth=1.5pt,
  innertopmargin=6pt, innerbottommargin=6pt,
  innerleftmargin=10pt, innerrightmargin=10pt,
  skipabove=6pt, skipbelow=4pt
]{formulabox}

% Mini example box
\newmdenv[
  backgroundcolor=exgreenbg,
  linecolor=exgreen, linewidth=1pt,
  innertopmargin=5pt, innerbottommargin=5pt,
  innerleftmargin=10pt, innerrightmargin=10pt,
  skipabove=4pt, skipbelow=4pt
]{examplebox}

% Watch-out box
\newmdenv[
  backgroundcolor=warnbg,
  linecolor=warnorange, linewidth=1pt,
  innertopmargin=5pt, innerbottommargin=5pt,
  innerleftmargin=10pt, innerrightmargin=10pt,
  skipabove=4pt, skipbelow=4pt
]{watchoutbox}

% Skill section heading
\newcommand{\skillheading}[1]{%
  \vspace{0.4cm}
  {\large\textbf{\textcolor{skillblue}{#1}}}
  \vspace{0.1cm}
  \hrule height 1pt
  \vspace{0.2cm}
}

\akheader{Compositions and Symmetry}

\begin{document}
\aktitleblock{Compositions of Rigid Motions and Symmetry}{Geometry}{}
\input{qa_transformations_curr347}

\noindent\small\textbf{What this sheet is teaching, and what to listen for.} The sheet builds
one argument in three stages: \emph{do} a composition (1--3), \emph{name} the single
transformation it amounts to (4--6), and then use that machinery to describe a figure's own
symmetries (8--10). The commonest error is not arithmetic, it is order --- ``reflect then
translate'' and ``translate then reflect'' generally land in different places, and problem 7
asks the student about exactly that. When a description is wrong, check whether it names all
of kind, size and centre or direction; a rotation without a centre is not an
answer.\normalsize

\problem{\textbf{One point, two steps.} $A(3, 4)$, reflect in the $x$-axis then translate by
$\langle 2, 5 \rangle$.

\par\smallskip
Reflecting in the $x$-axis keeps the $x$-coordinate and negates the $y$-coordinate:
$A' = (3, -4)$.
\par
Translating adds the vector componentwise: $x$ goes to $3 + 2 = \ans{5}$ and $y$ goes to
$-4 + 5 = \ans{1}$, so $A'' = (5, 1)$.
\par
\textbf{Do one step at a time, and plot both.} A student who jumps straight to the final point
usually applies the translation to the original $A$ and lands on $(5, 9)$. Plotting $A'$ makes
that impossible to do by accident.}

\problem{\textbf{A segment through a composition.} $P(1,1)$, $Q(4,5)$; translate by
$\langle 3, -2 \rangle$, then reflect in the $y$-axis.

\par\smallskip
Translation: $P' = (4, -1)$ and $Q' = (7, 3)$.
\par
Reflection in the $y$-axis negates the $x$-coordinate: $P'' = (-4, -1)$ and $Q'' = (-7, 3)$.
\par
Original length: $PQ = \sqrt{(4-1)^2 + (5-1)^2} = \sqrt{9 + 16} = \ans{5}$.
\par
Image length: $P''Q'' = \sqrt{(-7-(-4))^2 + (3-(-1))^2} = \sqrt{9 + 16} = \ans{5}$.
\par
\textbf{The equal lengths are the point, not a coincidence.} Both steps are rigid motions, so
the composition is one too, and a rigid motion cannot change a distance. If a student's two
lengths differ, the error is in the image coordinates --- the distance formula has just caught
it for them, which is a habit worth building.}

\problem{\textbf{A triangle through a composition.} $A(-4,1)$, $B(-1,1)$, $C(-1,3)$; reflect in
the $y$-axis, then translate by $\langle 0, -4 \rangle$.

\par\smallskip
Reflection in the $y$-axis: $A' = (4, 1)$, $B' = (1, 1)$, $C' = (1, 3)$.
\par
Translation down $4$: $A'' = (4, -3)$, $B'' = (1, -3)$, $C'' = (1, -1)$.
\par
\textbf{(a)} $\overline{A''B''}$ is horizontal, so its length is just the gap in $x$:
$4 - 1 = \ans{3}$. That matches $AB$ on the original, as it must.
\par
\textbf{(b)} Midpoint of $\overline{A''C''}$: average the coordinates,
$\left( \frac{4+1}{2}, \frac{-3+(-1)}{2} \right) = \ans{(2.5,\, -2)}$.
\par
\textbf{A check that costs nothing.} The midpoint of $\overline{AC}$ on the original is
$(-2.5, 2)$. Reflecting that in the $y$-axis and dropping it $4$ gives $(2.5, -2)$ --- the same
point. Midpoints go where the transformation sends them, so a student can verify part (b)
without recomputing anything.}

\problem{\textbf{Two parallel mirrors.} Reflect $P(0,2)$ in $x = 1$, then in $x = 4$.

\par\smallskip
\textbf{(a)} $P$ is $1$ unit left of $x=1$, so its image is $1$ unit right of it:
$P' = (2, 2)$. Then $P'$ is $2$ units left of $x = 4$, so $P'' = (6, 2)$.
\par
\textbf{(b)} $P(0,2)$ to $P''(6,2)$ is a horizontal move of $\ans{6}$ units.
\par\smallskip
\textbf{(c)} (Open description --- listen for all three parts of it.)
\par
\textbf{A model answer.} ``A translation of $6$ units to the right --- the vector
$\langle 6, 0 \rangle$. Nothing is flipped: two reflections undo each other's mirror-image
effect, so the final figure faces the same way as the original.''
\par
\textbf{The general result behind it.} The mirrors are $3$ units apart and the point moved $6$
units --- exactly twice the gap, in the direction from the first mirror to the second. That is
true for every point and every pair of parallel mirrors, and it does not depend on where the
point started. Worth asking: what if the reflections were done in the other order?
(A translation of $6$ units to the \emph{left}.)
\par
\textbf{Marking guide.} Full credit needs the kind (translation), the size ($6$) and the
direction (right / positive $x$). ``A translation'' alone is half an answer.
\par
\textbf{If they answered $3$ in part (b):} they gave the distance between the mirror lines
instead of the distance the point travelled. Have them measure $P$ to $P''$ on the grid.}

\problem{\textbf{Two perpendicular mirrors.} Reflect $P(3,5)$ in the $x$-axis, then in the
$y$-axis.

\par\smallskip
\textbf{(a)} Reflecting in the $x$-axis gives $(3, -5)$; reflecting that in the $y$-axis gives
$P'' = (-3, -5)$.
\par
\textbf{(b)} Midpoint of $\overline{PP''}$:
$\left( \frac{3 + (-3)}{2}, \frac{5 + (-5)}{2} \right) = \ans{(0,\, 0)}$ --- the origin.
\par\smallskip
\textbf{(c)} (Open description --- listen for kind, angle and centre.)
\par
\textbf{A model answer.} ``A rotation of $180^\circ$ about the origin --- a half turn. Part (b)
shows why: the origin is exactly halfway between $P$ and its final image, and a half turn about
a point is precisely the motion that sends every point to the other side of that point, the
same distance away.''
\par
\textbf{Marking guide.} Full credit for naming the transformation, the angle (or ``half
turn''), and the centre. ``A rotation about the origin'' without the angle is incomplete ---
$90^\circ$ about the origin sends $P$ somewhere else entirely. ``Point reflection in the
origin'' is also full credit; it is the same motion under another name.
\par
\textbf{Why the direction does not need stating here.} A half turn clockwise and a half turn
counterclockwise land in the same place, so $180^\circ$ is the one rotation angle that needs no
direction. Problem 6 does need one.}

\problem{\textbf{Mirrors that meet at $45^\circ$.} $\triangle RST$ with $R(1,1)$, $S(5,1)$,
$T(5,4)$; reflect in $y = x$, then in the $x$-axis.

\par\smallskip
Reflecting in $y = x$ swaps the coordinates: $R'(1,1)$, $S'(1,5)$, $T'(4,5)$.
\par
Reflecting that in the $x$-axis negates the $y$-coordinates: $R''(1,-1)$, $S''(1,-5)$,
$T''(4,-5)$.
\par
\textbf{(a)} $\triangle RST$ has a horizontal leg of $4$ and a vertical leg of $3$, so its area
is $\tfrac{1}{2}(4)(3) = \ans{6}$ square units.
\par
\textbf{(b)} The image has a vertical leg of $4$ and a horizontal leg of $3$, so its area is
$\tfrac{1}{2}(4)(3) = \ans{6}$ square units --- unchanged, as it must be.
\par\smallskip
\textbf{(c)} (Open description --- listen for kind, angle, direction and centre.)
\par
\textbf{A model answer.} ``A rotation of $90^\circ$ clockwise about the origin. In coordinates
the two steps send $(a, b)$ to $(b, a)$ and then to $(b, -a)$, and $(a,b) \mapsto (b,-a)$ is
the rule for a quarter turn clockwise about the origin. Checking on $T$: $(5,4)$ becomes
$(4,-5)$, which is what the drawing shows.''
\par
\textbf{Marking guide.} Full credit requires the angle \emph{and} the direction. $90^\circ$
counterclockwise is the composition done in the other order, so it is a genuinely wrong answer
rather than a slip. ``$270^\circ$ counterclockwise'' is the same motion and earns full credit.
\par
\textbf{The angle between the mirrors is the tell.} The lines $y = x$ and the $x$-axis meet at
$45^\circ$, and the resulting rotation is $90^\circ$ --- twice the angle between them, about
their intersection point. Problem 5 is the same rule with mirrors at $90^\circ$ giving a
$180^\circ$ rotation, and problem 4 is the limiting case with mirrors that never meet.}

\problem{\textbf{Translation, then a half turn.} $M(2,3)$, translate by
$\langle 5, -1 \rangle$, then rotate $180^\circ$ about the origin.

\par\smallskip
Translation: $M' = (2 + 5,\ 3 - 1) = (7, 2)$.
\par
A half turn about the origin negates both coordinates: $x$ goes to $-(2 + 5) = \ans{-7}$ and
$y$ goes to $-(3 + (-1)) = \ans{-2}$, so $M'' = (-7, -2)$.
\par
\textbf{Order matters, and here is the proof.} Doing the half turn first sends $M(2,3)$ to
$(-2,-3)$, and translating that by $\langle 5, -1 \rangle$ gives $(3, -4)$ --- a different
point from $(-7,-2)$. So the answer to the small question at the bottom of the problem is
\emph{no}. Composition is not commutative, and this is the cheapest example to show it.
\par
\textbf{A rotation of $180^\circ$ about the origin is worth memorising as a rule:}
$(x, y) \mapsto (-x, -y)$. Students who rotate by tracing on the grid get there too, but the
rule is what makes problem 5's description available.}

\problem{\textbf{Symmetry of a regular hexagon.}

\par\smallskip
\textbf{(a)} Two families of lines: three run from a vertex to the opposite vertex, and three
run from an edge midpoint to the opposite edge midpoint. That is $2 \times 3 = \ans{6}$ lines
of symmetry.
\par
\textbf{(b)} Six vertices sit at equal spacings around the centre, so turning by one sixth of a
full turn lands each vertex on the next one:
$360 \div 6 = \ans{60\text{ degrees}}$.
\par
\textbf{The general rule.} A regular $n$-gon has exactly $n$ lines of symmetry and rotational
symmetry of order $n$, with smallest angle $360^\circ/n$. Both parts of this problem are that
one rule with $n = 6$, which is worth saying out loud so the student is not memorising the
hexagon.
\par
\textbf{Why the two families look different but count the same.} With an even number of sides,
a vertex-to-vertex line has a vertex at each end and an edge-to-edge line has an edge midpoint
at each end. With an odd number of sides every line of symmetry runs vertex to edge midpoint,
so there is only one family --- and still exactly $n$ of them.
\par
\textbf{If they answered $180$ in part (b):} they found \emph{a} rotation that works rather
than the \emph{smallest} one. A half turn does map the hexagon onto itself; so do $120^\circ$,
$240^\circ$ and $300^\circ$. The question asks for the smallest, and every other one is a
multiple of it.}

\problem{\textbf{Symmetry of a parallelogram.} $E(0,0)$, $F(4,0)$, $G(6,3)$, $H(2,3)$.

\par\smallskip
\textbf{(a)} Midpoint of $\overline{EG}$:
$\left( \frac{0+6}{2}, \frac{0+3}{2} \right) = \ans{(3,\, 1.5)}$.
\par
\textbf{(b)} Midpoint of $\overline{FH}$:
$\left( \frac{4+2}{2}, \frac{0+3}{2} \right) = \ans{(3,\, 1.5)}$ --- the same point.
\par\smallskip
\textbf{(c)} (Open explanation --- listen to the reason.)
\par
\textbf{A model answer.} ``Both diagonals have midpoint $(3, 1.5)$, so call that point $Z$.
A half turn about $Z$ sends $E$ to $G$ and $F$ to $H$, because $Z$ is halfway along each of
those segments --- and it sends $G$ to $E$ and $H$ to $F$ for the same reason. Every vertex
lands on a vertex, so the whole parallelogram maps onto itself: it has rotational symmetry of
order $2$. But it has no line of symmetry. The two slanted sides $\overline{FG}$ and
$\overline{EH}$ both lean to the right rather than mirroring each other, so folding along a
vertical line through $Z$ does not bring them together, and folding along either diagonal sends
$E$ off the figure entirely.''
\par
\textbf{Marking guide.} Full credit needs the shared midpoint used as the centre of the half
turn, plus a \emph{reason} for the missing line symmetry --- either the leaning sides or a named
vertex whose mirror image is not a vertex of $EFGH$. Partial credit for the rotation argument
alone. Not sufficient: ``parallelograms don't have line symmetry'', which is a fact recalled
rather than a reason, and is in any case false for the rectangle and the rhombus.
\par
\textbf{A good follow-up.} Ask which parallelograms \emph{do} have line symmetry. A rectangle
has two, a rhombus has two, and a square has four --- and all of them still have the half turn,
because the diagonals of every parallelogram bisect each other.}

\problem{\textbf{Why perpendicular mirrors make a half turn.} Lines $x = 2$ and $y = 3$,
meeting at $C(2,3)$; $K = (4,1)$.

\par\smallskip
\textbf{(a)} $K$ is $2$ units right of $x = 2$, so reflecting puts it $2$ units left:
$K' = (0, 1)$. Then $K'$ is $2$ units below $y = 3$, so reflecting puts it $2$ units above:
$K'' = (0, 5)$.
\par
\textbf{(b)} Midpoint of $\overline{KK''}$:
$\left( \frac{4+0}{2}, \frac{1+5}{2} \right) = \ans{(2,\, 3)}$ --- exactly the crossing point
$C$.
\par\smallskip
\textbf{(c)} (Open justification --- the general argument is what is being marked.)
\par
\textbf{A model answer.} ``Put the crossing point at the origin, so the two mirrors are the
$y$-axis and the $x$-axis and a general point is $(x, y)$. Reflecting in the vertical mirror
reverses the horizontal displacement and leaves the vertical one alone: $(x, y) \to (-x, y)$.
Reflecting that in the horizontal mirror does the opposite: $(-x, y) \to (-x, -y)$. So after
both reflections \emph{both} displacements from the crossing point have been reversed, and
$(x, y) \mapsto (-x, -y)$ is exactly a $180^\circ$ rotation about that point. Nothing in the
argument used particular numbers, so it holds for every point, and therefore for every figure.''
\par
\textbf{Marking guide.} The question asks for an argument that covers \emph{every} point, so
re-checking $K$ is not sufficient on its own --- that is part (a). Full credit for the
coordinate argument above, or for a clear description in words of the two displacements
reversing one after the other, or for the midpoint observation \emph{provided} the student says
why the crossing point is the midpoint for every point and not just this one. Partial credit
for a correct statement of the result with a single worked instance as the only support.
\par
\textbf{If the student is stuck:} ask them to redo part (a) with a second point of their own
choosing, then a third, and to say what stayed the same each time. The pattern --- both
displacements from $C$ flip sign --- is the argument, and most students can state it once they
have seen it happen twice.}

\end{document}
